% =====================================================================
%  capitulo-nuevo.tex — Plantilla para un capítulo de la documentación
%  ---------------------------------------------------------------------
%  1. Copiar este archivo a capitulos/NN_mi-capitulo.tex
%  2. Cambiar título, etiqueta y contenido.
%  3. Registrarlo en main.tex con \input{capitulos/NN_mi-capitulo}
%  4. Crear su carpeta de imágenes: figuras/NN_mi_capitulo/
%  5. Compilar con: make
%
%  NO cargar paquetes aquí: si hace falta uno, va en estilos/cisa.sty.
% =====================================================================
\chapter{Título del capítulo}
\label{cap:mi-capitulo}

\section{Primera sección}

Texto introductorio. Explicar el problema antes de la solución.

\subsection{Un detalle}

Para rutas y archivos usar \path{ml/hilos.py}. Para código corto,
\texttt{N\_VOTOS} (con el guion bajo escapado).

\begin{nota}
Las cajas \texttt{nota}, \texttt{aviso} y \texttt{peligro} sirven para
resaltar información que no debe perderse entre el texto.
\end{nota}

\section{Una figura}

\begin{figure}[htbp]
  \centering
  % \includegraphics[width=0.8\linewidth]{NN_mi_capitulo/mi-imagen.png}
  \caption{Pie de figura que explica qué se ve.}
  \label{fig:mi-figura}
\end{figure}

Como muestra la figura~\ref{fig:mi-figura}, \ldots

\section{Una tabla}

\begin{center}
\begin{tabular}{@{}ll@{}}
  \toprule
  \textbf{Columna} & \textbf{Qué significa} \\
  \midrule
  Dato 1 & Explicación. \\
  Dato 2 & Explicación. \\
  \bottomrule
\end{tabular}
\end{center}

\section{Un bloque de código}

\begin{lstlisting}[language=bash, caption={Qué hace este comando},
                   label={lst:mi-comando}]
docker compose up -d
\end{lstlisting}

\section{Cierre}

Resumir la idea principal y enlazar con el capítulo siguiente
(capítulo~\ref{cap:referencia}).
